\documentclass[11pt]{article}
\usepackage[margin=1in]{geometry}
\usepackage{amsmath,amssymb,amsthm,hyperref}
\newtheorem{theorem}{Theorem}\newtheorem{lemma}[theorem]{Lemma}
\newtheorem{corollary}[theorem]{Corollary}
\newcommand{\E}{\mathbb E}\newcommand{\R}{\mathbb R}
\newcommand{\relu}{\operatorname{ReLU}}
\title{Locality of parity transfer in a learned neural representation}
\author{ChatGPT-6 Astra\thanks{Independent AI research. Sole author. Verification scope and provenance are stated in the accompanying research collection.}}
\date{18 November 2026; independent research draft}
\hypersetup{pdftitle={Locality of parity transfer in a learned neural representation},pdfauthor={ChatGPT-6 Astra}}
\emergencystretch=3em
\begin{document}\maketitle
\begin{abstract}
A representation can transfer efficiently to neighboring tasks while
remaining expensive for other tasks of equally small neural complexity.
We make this distinction explicit for a shallow ReLU network pretrained
on full parity. The accompanying cold-Gaussian posterior theorem supplies
polynomial-norm predictors for every parity obtained by deleting a fixed
even number of coordinates. Conversely, permutation symmetry and a trace
budget imply an exponential RKHS approximation-norm lower bound for any
specified parity on about half the coordinates, with high probability over
pretraining. This negative result holds at finite width and needs no
condensation limit. A fresh Gaussian linear readout with fixed positive
likelihood variance then requires exponentially many labels for small
squared prediction error on that distant task. An averaged sequential
version of the same readout learns the neighboring tasks with polynomially
many labels. The lower bound concerns frozen-feature regression, not
end-to-end classification mistakes.
\end{abstract}

\section{A general trace and symmetry principle}
Let $X=\{-1,1\}^r$ carry uniform probability and let
$\chi_T(x)=\prod_{i\in T}x_i$. Kernels act on $L^2(X)$ with its
probability normalization. For a positive semidefinite kernel $K$, put
\[
 \lambda_T(K)=\langle\chi_T,K\chi_T\rangle,
 \qquad \operatorname{tr}K=\E_xK(x,x).
\]
Every RKHS predictor satisfies
\[
 |\langle g,\chi_T\rangle|
       \le \|g\|_{\mathcal H(K)}\sqrt{\lambda_T(K)}. \tag{1}
\]
This follows by writing $g$ in feature coordinates and applying
Cauchy--Schwarz; equivalently use the adjoint of the RKHS embedding.
Consequently $\|g-\chi_T\|_2\le a<1$ implies
\[
 \|g\|_{\mathcal H(K)}^2\ge(1-a)^2/\lambda_T(K),       \tag{2}
\]
with the natural interpretation when $\lambda_T=0$.

\begin{theorem}[A fixed target in a large symmetric family]\label{thm:trace}
Let $K$ be random and let $G$ be an event such that the joint law
of $(K,\mathbf1_G)$ is invariant under permutations of the input
coordinates. Suppose $\operatorname{tr}K\le Q$ on $G$.
For every fixed $T$ of size $j$, every $0<\xi<1$, and $0\le a<1$,
\[
 \Pr\left(G\ \text{and}\
 \inf_{\|g-\chi_T\|_2\le a}\|g\|_{\mathcal H(K)}^2
              <\frac{(1-a)^2\xi\binom rj}{Q}\right)\le\xi.
 \tag{3}
\]
The infimum is over the RKHS of the realized kernel. The target is
fixed before that realization; no adversarial target selection is needed.
\end{theorem}
\begin{proof}
Bessel's inequality for the orthonormal Walsh family gives
$\sum_{|S|=j}\lambda_S(K)\le\operatorname{tr}K$.
Permutation invariance makes
$\E[\mathbf1_G\lambda_S(K)]$ the same for every such $S$.
Thus $\E[\mathbf1_G\lambda_T(K)]\le Q/\binom rj$.
Markov's inequality bounds the probability on $G$ that
$\lambda_T>Q/(\xi\binom rj)$ by $\xi$. Apply (2).
The proof also gives a deterministic fraction statement: on every
realization in $G$, at least $1-\xi$ of these characters have the
norm lower bound in (3).
\end{proof}

The argument is elementary and applies to any feature-learning method
with the stated symmetry and trace budget. It does not say that a
preselected arbitrary task is hard for every learned representation.
Symmetry is what turns a large hard fraction into a fixed-target
probability statement.

\section{Application to finite neural posteriors}
Consider
\[
 F(x)=\frac1N\sum_i a_i\relu(U_i\cdot(1,x)),\qquad
 E=\frac1{2N}\sum_i(a_i^2+\|U_i\|^2).
\]
Let $K$ be either the original frozen incoming kernel
$N^{-1}\sum h_{U_i}\otimes h_{U_i}$ or the even kernel with
$\psi_i(x)=[h_{U_i}(x)+h_{U_i}(-x)]/2$ in place of $h_{U_i}$.
Both satisfy the sharper average, rather than uniform-radius, bound
\[
 \operatorname{tr}K\le\frac1N\sum_i\|U_i\|^2\le2E.   \tag{4}
\]
Indeed $\E_x(U\cdot(1,x))^2=\|U\|^2$, ReLU does not increase
absolute value, and even projection is an $L^2$ contraction.

Pretrain on the full-parity target $\chi_{[r]}$ from uniform IID
inputs. A product Gaussian prior with common coordinate variance,
together with an observed-margin or square likelihood, is equivariant
under coordinate permutations. Its joint sample--posterior law therefore
has the symmetry in Theorem~\ref{thm:trace}.
This is a joint-law statement: a particular asymmetric training sample
need not have a permutation-invariant conditional posterior.

The cold-Gaussian hard-margin companion proves, for fixed even $r\ge6$,
margin $s>0$, and arbitrarily prescribed sufficiently small
$\epsilon>0$, explicit polynomial thresholds on prior scale, sample
size, and width with these properties:
\begin{enumerate}
\item On a common sample event of probability at least $1-\zeta$,
      the posterior assigns mass at least $1-\beta^N$ to a
      permutation-invariant good event $G$.
\item On $G$, $E\le B=(1+\epsilon)s/\rho$, where
      $\rho=\binom{r-2}{(r-2)/2}/(2^{r-1}\sqrt r)=\Theta(r^{-1})$.
\item For each fixed $k$, every degree-$(r-2k)$ target has a comparator
      in the even-feature RKHS with squared norm
      $B_k^2=O_k(r^{2k+1})$, squared feature radius $R^2=O(r^2)$,
      and any prescribed inverse-polynomial $L^2$ approximation error,
      by choosing the pretraining accuracy correspondingly.
\end{enumerate}
Here $G$ can be taken as the population hinge-loss and energy event;
these conditions themselves are invariant. The sample event need not
be separately included in $G$. The square-likelihood companion supplies
an alternative using a smooth likelihood and a population objective gap.
All these conclusions concern actual finite-width posteriors.

\begin{corollary}[Near-task and distant-task separation]\label{cor:separation}
Use a pretraining regime just described and let
$j=2\lfloor r/4\rfloor$. Fix in advance a set $T$ of size $j$.
With joint pretraining probability at least
$1-\zeta-\beta^N-1/r$, all the neighboring-task comparators hold and
\[
 \inf_{\|g-\chi_T\|_2\le1/4}\|g\|_{\mathcal H(K_e)}^2
               \ge\frac{9\binom rj}{32Br}
               \ge\frac{9\,2^r}{64Br(r+1)}.            \tag{5}
\]
The lower bound is $\Omega(2^r/r^3)$ at fixed $s$. It also holds
for the original unsymmetrized frozen kernel.
\end{corollary}
\begin{proof}
Apply (3) with $Q=2B$, $a=1/4$, and $\xi=1/r$.
The chosen binomial coefficient is central or one step from central,
so it is at least $2^{r-1}/(r+1)$.
Intersect with $G$, whose failure probability is at most
$\zeta+\beta^N$.
For a simultaneous statement about both kernels one need not union
bound: their quadratic forms on every even character are identical,
because even projection leaves that character unchanged.
\end{proof}

The source, neighboring targets, and distant target all have shallow
ReLU teachers of width $O(r)$ and polynomially bounded coefficients.
For a parity on $j$ variables, interpolate its alternating values at
$-j,-j+2,\ldots,j$ with the hinge basis of the coordinate sum.
Thus the separation is in transfer geometry, not the existence of a
small target network. The distant target can in fact have fewer
relevant coordinates than the source.

\section{An exponential downstream regression obstruction}
For either realized frozen feature map, use its normalized coordinates
with a standard Gaussian readout prior and Gaussian likelihood variance
$v>0$. For $n$ labels $y_i\in\{-1,1\}$ the posterior mean $g_n$
minimizes
\[
 \sum_{i=1}^n(y_i-g(x_i))^2+v\|g\|_{\mathcal H(K)}^2.
\]
Comparison with zero gives $\|g_n\|_{\mathcal H(K)}^2\le n/v$.
This holds for every dataset, including adaptive choices of inputs.
Therefore on (5), whenever
\[
                 n<\frac{9v\,2^r}{64Br(r+1)},           \tag{6}
\]
its uniform mean squared error is greater than $1/16$.
Expected squared error of a posterior prediction draw is at least the
error of its mean. Any convex average of posterior means from prefixes
of that dataset has the same norm ceiling and the same obstruction.

If downstream inputs are uniform and independent of the past, the
bound holds conditionally on every past dataset. Hence for every horizon
below (6), the expected cumulative squared prediction loss of the
sequential posterior mean exceeds one sixteenth of the horizon.
This is a squared-loss statement. Small amplitude with the correct
sign can evade an approximation-norm lower bound, so no analogous
zero--one mistake conclusion follows from this proof alone.

\section{Polynomial neighboring-task learning with the same readout}
An elementary sequential-ridge identity supplies a useful positive
comparison. Let a fixed feature map have squared radius at most $R^2$,
and let $g_t$ be the Gaussian posterior mean after $t$ labels, with the
same variance $v$. For every deterministic sequence and comparator
$q\in\mathcal H(K)$,
\[
 \sum_{t=1}^n(y_t-g_{t-1}(x_t))^2
 \le(1+R^2/v)\left\{\sum_{t=1}^n(y_t-q(x_t))^2
                                      +v\|q\|_{\mathcal H(K)}^2\right\}.
 \tag{7}
\]
To verify this, write $A_t=vI+\sum_{i\le t}\phi_i\phi_i^\top$,
$b_t=\sum_{i\le t}y_i\phi_i$, and
$M_t=\min_w\{\sum_{i\le t}(y_i-w\cdot\phi_i)^2+v\|w\|^2\}$.
The rank-one inverse formula gives
\[
 M_t-M_{t-1}
 =\frac{(y_t-g_{t-1}(x_t))^2}
        {1+\phi_t^\top A_{t-1}^{-1}\phi_t}.
\]
The denominator is at most $1+R^2/v$. Sum this equality and compare
$M_n$ with the objective at $q$ to obtain (7).
The feature dimension can be arbitrary finite width.

For IID downstream samples labeled by a fixed target $f$, define
$\bar g_n=n^{-1}\sum_{t=1}^n g_{t-1}$ as a function on the input
space. Conditional independence, (7), and convexity give
\[
 \E\|\bar g_n-f\|_2^2
 \le(1+R^2/v)
       \left\{\|q-f\|_2^2+v\|q\|_{\mathcal H(K)}^2/n\right\}.
 \tag{8}
\]
The expectation here is only over the downstream data; the pretrained
features are fixed. This estimator is explicitly an average of Bayesian
posterior means, rather than a claim about monotonicity of the final
posterior mean's risk.

For a neighboring parity, choose the pretraining accuracy so that
its comparator error is at most
$\delta/\sqrt{2(1+R^2/v)}$ and take
\[
 n\ge 2(v+R^2)B_k^2/\delta^2.
\]
Then (8) is at most $\delta^2$, with
$n=O_k(r^{2k+3}\delta^{-2})$ at fixed $v$.
The extra inverse-polynomial pretraining accuracy is permitted by the
finite cold-prior theorem, so pretraining remains polynomial as well.
This same averaged estimator remains subject to the distant-task lower
bound (6). Thresholding gives a positive classification guarantee for
the neighboring tasks, while the negative distant-task assertion remains
one about squared prediction error.

\section*{Scope and verification}
The trace, symmetry, and sequential-ridge arguments are elementary
consequences of orthogonality and regularized least squares; they are
not presented as new general tools. Their role is to give a finite-width
local-versus-distant transfer separation for the actual neural posterior
constructed in the companions. No efficient pretraining sampler is
claimed. The prior scale and the width/sample relation are part of that
model. Fine-tuning hidden weights, choosing a very diffuse new readout
prior, changing downstream likelihood variance with dimension, or using
an algorithm outside the stated norm budget can evade this conclusion.
These are conventional same-session proofs without independent review.
The accompanying Lean file checks finite orthogonal-target capacity,
coefficient Cauchy--Schwarz, and the ridge norm budget. The posterior
symmetry argument, its probabilistic application, and the sequential
readout identity are not formalized there.
\end{document}
